\documentclass[12pt]{article}
\usepackage[utf8]{inputenc}
\usepackage[T1]{fontenc}
\usepackage{hyperref}
\title{us-Saare Language Resources for K-12 Curriculum Development}
\author{Compiled by Claude Sonnet 5 (claude-sonnet-5) for Andr\'as Kornai}
\date{}
\begin{document}
\maketitle

\section{Introduction}
us-Saare (ISO 639-3: uss) is the western variety of Hun-Saare, a Kainji language of Kebbi and Niger States, Nigeria; the eastern variety is Ut-Hun (uth). Speakers call both Saare. A figure of about 73,000 (1985) is cited, also associated with the Dukawa people. Only general and mission-organization pages were found; no educational materials.

\section{Primary Educational Resources}
None found.

\section{Linguistic and Reference Materials}
\begin{itemize}
\item Wikipedia, Hun-Saare language: \url{https://en.wikipedia.org/wiki/Hun-Saare_language}
\end{itemize}
No dictionary, grammar or orthography was found in search results.

\section{Religious and Community Resources}
\begin{itemize}
\item Joshua Project language page (lists alternate names Western Duka/Dukawa and dialects Dugge, Dukku, Iri, Sandiri, Tungan Bunu, Upoki): \url{https://www.joshuaproject.net/index.php/languages/uss}
\end{itemize}

\section{Digital and Audio Resources}
Global Recordings Network pages for dialects (contents not inspected):
\begin{itemize}
\item Tungan Bunu: \url{https://globalrecordings.net/en/language/24008}
\item Dugge: \url{https://globalrecordings.net/en/language/24004}
\item Upoki: \url{https://globalrecordings.net/en/language/24009}
\end{itemize}

\section{Curriculum Development Context}
The language appears to have a sizeable speaker base but little documented educational material online. Curriculum work would need primary documentation with the community.
\end{document}